\section{Lawfulness on the audited shell}\label{sec:lawfulness}

We now state the first closed theoremlet in the paper: the continuous full-loop lawfulness theoremlet on the audited shell. This result is the base structural theorem in the package. It does not yet concern strict theory extension or thermodynamic consequence. Its role is to show that the shell-stable class supports a well-defined, nondegenerate, six-mechanism update law whose dynamics do not collapse immediately into a trivial frozen regime.

\subsection{Lawful shell data}

Recall from Section~\ref{sec:preliminaries} that a substrate state is
\[
x=(K,\ell,\mathfrak{p},\tau,b,\eta)\in \mathcal{S}_{\mathrm{aud}},
\]
with continuous operator/kernel state \(K\), active lens \(\ell\), active packaging \(\mathfrak{p}\), timescale \(\tau\), budget variable \(b\), and auxiliary selector/hysteresis state \(\eta\). The update law
\[
F:\mathcal{S}_{\mathrm{aud}}\to \mathcal{S}_{\mathrm{aud}}
\]
is generated by the six interacting mechanisms \(P1\)--\(P6\), viewed here as a deterministic hybrid update with switching and hysteresis on the audited shell \cite{GoebelSanfeliceTeel2012Hybrid}. The shell-stable class is the class denoted internally by
\[
\texttt{continuous\_full\_loop\_lawful\_kernel\_class\_shell\allowbreak\_stable}.
\]

For manuscript purposes, the hypotheses that matter here are the following:
\begin{enumerate}
    \item the operator state remains in a controlled shell of kernel states,
    \item all six mechanisms remain causally active on that shell,
    \item lens and packaging selectors remain nondegenerate under hysteresis,
    \item the budget and timescale variables remain in a lawful range,
    \item the resulting dynamics exhibit persistent variation and do not undergo fast collapse.
\end{enumerate}
These hypotheses are exactly the shell-level assumptions frozen in the internal theorem package; no broader class is claimed in this section.

\begin{theorem}[Continuous full-loop lawfulness on the audited shell]\label{thm:lawfulness}
Let \(\mathcal{S}_{\mathrm{aud}}\) be the audited shell-stable class of continuous substrate states described above. Assume that the full update law \(F\) satisfies the shell-stable hypotheses just listed: well-posed continuous operator dynamics, full six-mechanism causal closure, hysteresis-stable lens and packaging selection, lawful budget/timescale evolution, and absence of fast collapse on the shell. Then the following hold on \(\mathcal{S}_{\mathrm{aud}}\).

\begin{enumerate}
    \item The update law is well-defined and remains inside the shell.
    \item Lens competition and packaging competition remain nondegenerate.
    \item Budget and timescale dynamics remain active and lawful.
    \item All six mechanisms \(P1\)--\(P6\) remain causally active in the resulting dynamics.
    \item The shell supports a nontrivial lawful exploratory regime rather than immediate collapse into a degenerate frozen state.
\end{enumerate}

In particular, \(\mathcal{S}_{\mathrm{aud}}\) supports a continuous full-loop lawful regime.
\end{theorem}

\begin{proof}[Proof sketch]
We indicate the structure of the closure; implementation-level estimates are deferred to the supporting package.

\smallskip
\noindent\emph{Step 1: well-posed state space.}
The state variables \((K,\ell,\mathfrak{p},\tau,b,\eta)\) define a well-posed update state on the audited shell. The kernel state evolves in a controlled continuous class, while the selector, budget, and timescale variables live in corresponding bounded shell ranges. This is the well-posedness ingredient behind the state-space lemma of the internal package.

\smallskip
\noindent\emph{Step 2: refresh and selector consistency.}
The refresh/invalidation mechanism is consistent on the shell, so lens and packaging candidates can be recomputed without ambiguity. The hysteretic selectors therefore define stable lens and packaging dynamics rather than arbitrary switching. In particular, lens and packaging competition remain meaningful rather than collapsing into undefined or pathological selection.

\smallskip
\noindent\emph{Step 3: six-mechanism closure.}
The shell hypotheses ensure that each of the six mechanisms enters the update law essentially:
\(P1\) changes the operator state,
\(P2\) changes admissibility/support,
\(P3\) changes timescale/protocol state,
\(P4\) changes the active lens,
\(P5\) changes the active packaging and feeds into other update components,
and \(P6\) changes the budget/ledger state.
Each mechanism is therefore essential: the lawful loop is a genuine six-mechanism loop, not a reduced loop with redundant primitives.

\smallskip
\noindent\emph{Step 4: forward shell preservation.}
The shell assumptions imply that one step of the update law keeps the dynamics inside the lawful shell. In particular, kernel variation, selector margins, budget range, and timescale range remain in the admissible region. This is the shell-preservation ingredient later reused by the theorem package.

\smallskip
\noindent\emph{Step 5: noncollapse.}
Because selector competition, budget dynamics, and timescale adaptation remain active inside the shell, the resulting dynamics do not collapse into a trivial frozen regime. The audited shell therefore carries a nontrivial lawful exploratory regime rather than a transient simulation artifact.

Combining these points proves the theoremlet.
\end{proof}

\begin{figure}[t]
\centering
\small
\begin{tabular}{@{}l r@{}}
\toprule
\multicolumn{2}{@{}l}{\textbf{Continuous substrate regime on the audited shell}} \\
\midrule
Base witness kernel dimension        & 16 \\
Shell witness kernel dimension       & 20 \\
Base witness lens switches           & 249 \\
Base witness packaging switches      & 249 \\
Base witness budget range            & $[2.714,\;12.0]$ \\
Base witness $\tau$ range            & $[0.600,\;0.982]$ \\
Shell witness budget range           & $[2.739,\;12.0]$ \\
Shell witness $\tau$ range           & $[0.600,\;0.861]$ \\
Shell exits on both witnesses        & 0 \\
All six primitives active            & yes \\
\bottomrule
\end{tabular}
\caption{Continuous full-loop regime on the audited shell.
The two frozen witnesses remain inside the shell-stable class, exhibit persistent
variation, nondegenerate lens and packaging competition, active
budget/timescale dynamics, and no shell exits.
This figure supports only the continuous full-loop lawfulness
theoremlet on the audited shell.}
\label{fig:full-loop-regime}
\end{figure}

\begin{figure}[t]
\centering
\small
\begin{tabular}{@{}l c c l@{}}
\toprule
\multicolumn{4}{@{}l}{\textbf{Primitive knockout comparison on the audited shell}} \\
\midrule
Removed & Lens sw. & Pkg.\ sw. & Main degradation \\
\midrule
none & 249 & 249 & reference lawful loop \\
P1   & 249 & 249 & $\tau$ switching collapse \\
P2   & 249 & 249 & stabilization change and damping \\
P3   &   0 &   0 & selector and $\tau$ collapse \\
P4   &   0 & 249 & lens collapse \\
P5   & 249 &   0 & packaging, budget, and $\tau$ degradation \\
P6   & 249 & 249 & budget collapse \\
\bottomrule
\end{tabular}
\caption{Leave-one-primitive-out comparison on the audited shell.
The full loop remains lawful, while removal of any one primitive
materially degrades the regime.
In particular, removing $P4$ collapses lens switching,
removing $P5$ collapses packaging competition and narrows the budget range,
and removing $P6$ collapses budget dynamics.
This figure supports only the six-primitive closure component
of the lawfulness theoremlet.}
\label{fig:primitive-knockout-closure}
\end{figure}


Theorem~\ref{thm:lawfulness} supplies the base dynamical regime on which the rest of the paper is built. The cocycle-pressure, strict-extension, and conditional-disintegration theoremlets all presuppose that the shell carries a lawful full loop; without it, they would sit on an ill-defined substrate.
